\begin{tikzpicture}[
        xscale=1.75, 
        yscale=0.9, 
        >=Stealth,
        % 样式定义
        survivor/.style={circle, fill=blue!70!black, inner sep=1.2pt},
        other-node/.style={circle, draw=blue!70!black, fill=white, inner sep=1pt},
        killed-node/.style={circle, draw=blue!70!black, fill=white, inner sep=1pt, cross out, thick, minimum size=3pt},
        edge-live/.style={thick, blue!70!black},
        edge-killed/.style={thick, dashed, blue!70!black}, 
        tube-fill/.style={blue!10, opacity=0.4},
        tube-edge/.style={thin, black},
        ref-dash/.style={thin, gray!60, dashed}
    ]

        % 1. 坐标轴
        \draw[->, thick] (0, -2.5) -- (0, 3.8) node[above] {$\mathcal{V}(u)$};
        \draw[->, thick, dashed] (0, 0) -- (6.5, 0) node[right] {$n$};

        % 2. 阶段 I: 水平管道 [0, 1.5]
        \fill[tube-fill] (0, -2.0) rectangle (1.5, 0);
        \draw[tube-edge] (0, -2.0) -- (1.5, -2.0) node[pos=0, left, black] {$-b \log N$};
        \draw[tube-edge] (0, 0) -- (1.5, 0);

        % 3. 阶段 II: 倾斜管道 [1.5, 5.0]
        \def\ystart{-1.8} 
        \def\ylow{1.2}  % 终点下界
        \def\ywin{1.8}  % 终点窗口界
        \def\yhigh{2.8} % 终点上界
        \fill[tube-fill] (1.5, \ystart) -- (5.0, \ylow) -- (5.0, \yhigh) -- (1.5, \ystart+1.6) -- cycle;
        \draw[tube-edge] (1.5, \ystart) -- (5.0, \ylow); 
        \draw[tube-edge] (1.5, \ystart+1.6) -- (5.0, \yhigh); 
        
        \node[survivor] (y-node) at (1.5, \ystart) {};
        \node[left, xshift=-2pt] at (1.5, \ystart) {$y$};

        % 4. 均匀高频分支 (步长 0.3)
        \node[survivor] (n0) at (0, 0) {};

        % --- 阶段 I 路径 ---
        \foreach \i/\x/\y in {1/0.3/-0.4, 2/0.6/-0.1, 3/0.9/-0.8, 4/1.2/-1.4} {
            \coordinate (p\i) at (\x, \y);
            \node[survivor] at (p\i) {};
        }
        \draw[edge-live] (n0) -- (p1) -- (p2) -- (p3) -- (p4) -- (y-node);
        
        % 阶段 I 杀掉的分支 (管道外)
        \draw[edge-killed] (p1) -- (0.6, 0.4); \node[killed-node] at (0.6, 0.4) {};
        \draw[edge-killed] (p3) -- (1.2, -2.3); \node[killed-node] at (1.2, -2.3) {};
        \draw[edge-killed] (p4) -- (1.5, -1); \node[killed-node] at (1.5, -1) {};

        % --- 阶段 II 路径 (向上移动) ---
        \foreach \i/\x/\y in {1/1.8/-1.5, 2/2.1/-1.2, 3/2.4/-0.8, 4/2.7/-0.5, 5/3.0/-0.1, 6/3.3/0.3, 7/3.6/0.6, 8/3.9/0.9, 9/4.2/1.1, 10/4.5/1.3, 11/4.8/1.5} {
            \coordinate (q\i) at (\x, \y);
            \node[survivor] at (q\i) {};
        }
        \coordinate (m-node) at (5.0, 1.6);
        \draw[edge-live] (y-node) -- (q1) -- (q2) -- (q3) -- (q4) -- (q5) -- (q6) -- (q7) -- (q8) -- (q9) -- (q10) -- (q11) -- (m-node);
        \node[survivor] at (m-node) {};

        % 阶段 II 存活但非极小的分支 (管道内)
        \draw[edge-live] (q6) -- (3.6, 1.1) -- (3.9, 1.5) -- (4.2, 1.8) -- (4.5, 2.1) -- (4.8, 2.3);
        \foreach \x/\y in {3.6/1.1, 3.9/1.5, 4.2/1.8, 4.5/2.1, 4.8/2.3} \node[other-node] at (\x, \y) {};

        % 阶段 II 杀掉的分支 (越界)
        \draw[edge-killed] (q9) -- (4.5, 0.3); \node[killed-node] at (4.5, 0.3) {};
        \draw[edge-killed] (4.2, 1.8) -- (4.5, 2.8); \node[killed-node] at (4.5, 2.8) {};
        \draw[edge-killed] (4.8, 2.3) -- (5.0, 2.1); \node[killed-node] at (5.0, 2.1) {};

        % 5. 时间节点辅助线
        \draw[dashed, black!30] (1.5, -2.5) -- (1.5, 3.5) node[above, black] {$\ell$};
        \draw[dashed, black!30] (5.0, -2.5) -- (5.0, 3.5) node[above, black] {$\ell_N$};

        % 6. 标注窗口与跨度 (向右侧引出)
        
        % 左侧时刻 l 窗口
        \draw[ultra thick, blue!80!black] (1.5, -2.0) -- (1.5, -1.3);
        \draw[decorate, decoration={brace, amplitude=4pt, mirror}] (2, -2.0) -- (2, -1.3) 
            node[midway, right, xshift=2pt] {$v \log N$};

        % --- 右侧引出线与标注 ---
        % 引出辅助水平线
        \draw[ref-dash] (5.0, \ylow) -- (6.2, \ylow);
        \draw[ref-dash] (5.0, \ywin) -- (5.8, \ywin);
        \draw[ref-dash] (5.0, \yhigh) -- (6.2, \yhigh);
        \draw[ref-dash] (1.53, -2.0) -- (2, -2.0);
        \draw[ref-dash] (1.53, -1.3) -- (2, -1.3);

        % 标注 M^k (紧贴 L_N 垂线右侧)
        \node[right, blue!70!black] at (5, 1.6) {$\mathcal{M}^{\mathrm{k}}_{\ell_N}$};

        % 标注 vL^1/3 (第一层大括号)
        \draw[ultra thick, blue!80!black] (5, \ylow) -- (5, \ywin);
        \draw[decorate, decoration={brace, amplitude=4pt, mirror}] (5.6, \ylow) -- (5.6, \ywin) 
            node[midway, right, xshift=1pt] {$vL^{1/3}$};

        % 标注 cL^1/3 (第二层更远的大括号)
        \draw[decorate, decoration={brace, amplitude=5pt, mirror}] (6.3, \ylow) -- (6.3, \yhigh) 
            node[midway, right, xshift=2pt] {$cL^{1/3}$};

        % 7. 阈值线
        \draw[red!80!black, thick, dashed] (0, 3.2) -- (5.5, 3.2) node[right] {$\zeta' \log N$};

    \end{tikzpicture}
